Yes, a GitHub repository with no public license automatically implies "All Rights Reserved" by default under international copyright law.
## Key Facts About No-License Repositories

* Exclusive Ownership: The creator of the code retains full copyright ownership.
* No Rights to Modify: Others cannot legally copy, modify, distribute, or use your code in their own projects.
* No Commercial Use: Others cannot sell or monetize your unlicensed code.

## What GitHub's Terms of Service Allows
Even without a license, hosting code publicly on GitHub grants other users a few basic, automated rights so the platform can function:

* Viewing and Forking: Users can view and fork your repository through the GitHub interface.
* Cloning: Users can clone the repository to their local machines to view the code.

## The Problem with No License
If you leave a repository unlicensed, no one else can legally use your project, even if the code is completely public. If your goal is open-source collaboration, you must explicitly add a license file (like MIT or Apache 2.0).
If you want to update your repository, let me know if you need help with:

* Choosing the right open-source license (e.g., MIT, GPL, Apache)
* The steps to add a license file directly through GitHub
* How to handle code if you already forked an unlicensed project


